\chapter{引论}
\section{研究背景}
\subsection{全球及国内期权市场的发展现状}

自1973年芝加哥期权交易所（CBOE）正式挂牌交易标准化期权以来，全球金融衍生品市场经历了从单一品种向多层次、高复杂度体系的深刻演变。在国际成熟市场中，期权已不仅是基础的风险对冲工具，更演化为资产价格发现、尾部风险管理与策略博弈的核心载体。近年来，随着全球宏观经济不确定性加剧、地缘政治冲突频发，机构投资者对极端尾部风险的防范需求日益迫切，推动全球期权成交量持续创下历史新高。特别是以零日期权（0DTE）为代表的超短期品种在美股市场的大规模流行，不仅深刻改变了市场流动性的微观结构，更对定价模型在极短时间尺度下的瞬时响应能力、对极端波动率偏斜（Skew）的精准捕捉能力提出了更为严苛的理论挑战，传统常数波动率模型与线性随机波动率模型的局限性愈发凸显。

与此同时，中国期权市场作为全球衍生品版图中的重要新兴力量，正处于跨越式发展的关键窗口期。在场内市场层面，自上证50ETF期权2015年成功试点以来，国内已快速构建起覆盖沪深300、中证500、中证1000等核心指数，以及科创50、创业板指等特色指数的权益期权矩阵；商品期权领域则实现了从有色金属、黑色系产业链到农产品、新能源原材料的全品类覆盖，为实体企业提供了丰富的风险对冲工具，成为实体经济稳健经营的重要金融支撑。在场外（OTC）衍生品领域，以中证报价系统为核心纽带的场外期权市场稳步发展，个性化、定制化的风险解决方案日益丰富，进一步完善了多层次资本市场的风险管理功能。

\subsection{期权定价模型与计算方法的演进}

期权定价理论的演进，构成了现代金融工程学发展的核心逻辑主线。自 Black、Scholes 与 Merton 提出经典的 BSM 定价模型以来，金融资产价格遵循几何布朗运动、波动率为常数的核心假设，为衍生品定价提供了简洁且可解析的理论框架，奠定了现代衍生品市场的理论基石。然而，随着全球金融市场复杂性的持续提升，BSM 模型中常数波动率的假设与真实市场特征产生了系统性背离。大量实证研究表明，全球各类资产的收益率分布普遍存在显著的 “尖峰厚尾” 特征，期权市场隐含波动率层面则呈现出随行权价格与到期期限发生系统性变化的 “波动率微笑” 与 “波动率偏斜” 现象，这一现象在市场危机与极端行情中表现得尤为突出，也推动了波动率模型从确定性向随机性的范式转换。

在随机波动率模型的早期探索中，Heston 模型（亦称 1/2 模型）凭借其对波动率均值回归特性的优雅刻画，以及标的对数价格特征函数的仿射闭式表达，成为全球衍生品定价的标准基准。随后，针对利率与外汇市场的复杂波动动态，SABR 模型通过引入随机波动率项与波动率 - 标的价格的相关结构，进一步强化了对隐含波动率曲线形态的捕捉能力，成为场外衍生品定价的主流模型。然而，传统的 1/2 模型在处理极短期期权的隐含波动率陡峭偏斜、以及极端行情下波动率的爆发式跳变时，其线性的扩散项结构显露出生硬的理论局限性 —— 线性漂移项难以解释波动率在极端危机时刻的非线性加速均值回归特征，也无法充分刻画高波动市场中波动率的厚尾分布特性。

在此背景下，3/2 随机波动率模型逐渐进入学术与业界的核心研究视野。与 Heston 模型不同，3/2 模型假设瞬时方差过程遵循非线性的幂律扩散过程，其扩散项与瞬时方差的 3/2 次方成正比，漂移项则呈现出非线性的均值回归结构。这种非线性设定赋予了模型极强的灵活性，使其能够更灵敏地反映市场在极端压力下的波动聚集效应，更精准地拟合不同期限、不同行权价下的波动率微笑形态，从而在股票、外汇、加密货币等跨资产类别的定价实践中展现出显著优于传统仿射模型的拟合精度。

伴随着定价模型维度的增加与结构的复杂化，数值计算方法的演进同步推动着金融工程实践的边界拓展。早期的有限差分法（Finite Difference Method, FDM）通过离散化偏微分方程（PDE），为美式期权等路径依赖型衍生品提供了稳健的定价路径，但在多因素随机波动率模型下往往面临 “维数灾难” 的挑战，计算复杂度随模型维度呈指数级上升。蒙特卡洛模拟（Monte Carlo Simulation）方法以其强大的通用性，成为处理高维模型、路径依赖型奇异期权的主流工具，但其收敛速度较慢，计算成本随精度要求的提升而显著上升。为兼顾定价效率与精度，条件蒙特卡洛（Conditional Monte Carlo）通过引入方差缩减技术，将标的价格过程与波动率过程解耦，极大地优化了随机波动率模型的模拟效果与收敛速度；其中，逆高斯（Inverse Gaussian, IG）近似精确模拟、Gamma 近似精确模拟、Poisson-Gamma 近似精确步进、拟指数（Quasi-Exponential, QE）近似精确步进等方法，已在 Heston 等仿射模型中展现出优异的精度与效率平衡，但现有研究尚未将其系统引入 3/2 纯扩散模型的定价框架中。近年来，基于特征函数的傅里叶变换定价方法（Fourier Transform Method）结合快速傅里叶变换（FFT）、傅里叶余弦（COS）等数值算法，实现了复数域内期权价格的高速解析求解，为 3/2 模型等非仿射随机波动率模型的实时定价、全曲面校准奠定了坚实的技术基础；但现有文献仍缺乏 COS 方法与各类近似精确模拟方法在 3/2 模型下的系统性性能对标，尤其是极端市场参数下的全场景鲁棒性检验。

\section{研究问题}

尽管 3/2 模型在刻画市场极端波动、拟合波动率曲面方面具备显著的理论优势，但其非仿射的数学结构也为期权定价带来了严峻的数值实现挑战。现有研究针对 3/2 模型提出了多类期权定价方法，涵盖时间步进类蒙特卡洛模拟、终端精确 / 近似精确模拟、傅里叶频域定价三大范式，累计形成十一类主流实施方案，但现有研究与实践中仍存在四大核心待解决问题：

第一，现有研究对 3/2 模型下各类定价方法的性能评价，多局限于平稳市场环境下的理想化学术基准算例，缺乏对方法全场景数值稳定性的系统性检验。多数标称 “精确” 的定价方法，仅能在文献给定的基准参数下完成理论验证，在高波动率之波动率、强价格 - 波动率相关、极短到期期限等极端参数组合下，普遍存在特殊函数数值溢出、分布拟合失效、计算结果非数值化（nan/inf）等问题，方法的工程适用边界尚未得到清晰界定。

第二，针对逆高斯（IG）、Gamma 近似精确模拟、Poisson-Gamma、QE 近似精确步进这类在仿射模型中表现优异的近似精确方法，现有研究尚未完成其在 3/2 纯扩散模型下的系统性理论推导、数值实现与全场景性能检验，更未厘清其相较于传统离散格式、终端精确模拟方法的优势边界与适用场景，导致这类兼具精度与效率的方法在 3/2 模型的定价实践中缺乏可落地的理论支撑。

第三，真实极端市场环境下的定价方法实证再检验不足，理想化算例中的理论精度不等于真实市场中的实践性能。加密货币等极端波动市场的真实参数，往往超出学术基准算例的参数范围，极端行情下模型参数的结构性突变，会进一步放大不同定价方法的性能分化。现有研究尚未厘清不同方法在真实市场极端冲击下的定价精度、计算效率与鲁棒性，无法为极端行情下的衍生品定价与风险管理提供可靠的方法选择依据。

第四，3/2 模型在高波动市场中的实证适配性研究仍存在不足，缺乏极端事件全周期的模型参数动态演化分析。现有研究聚焦于 3/2 模型的静态参数校准与定价公式推导，尚未针对极端崩盘事件构建 “长期固定参数估计 - 时变参数横截面校准” 的完整分析框架，也未验证纯扩散框架下 3/2 模型对极端行情下波动率微笑形态演化的捕捉能力，模型在高波动市场中的实践价值尚未得到充分的实证支撑。

基于此，本文的核心研究问题聚焦于三个层面：其一，首次系统厘清 3/2 纯扩散模型下，包括 IG、Gamma、Poisson-Gamma、QE 近似精确方法与 COS 方法在内的十一类主流期权定价方法的理论逻辑、误差来源与收敛特性，通过多场景数值实验量化不同方法的定价精度、计算效率与全场景数值稳定性，形成清晰的方法性能分层体系；其二，基于真实极端市场事件校准的模型参数，完成定价方法的极端场景再检验，对比理想化算例与真实市场中方法性能的差异，明确不同方法的工程适用边界；其三，验证 3/2 模型在加密货币极端行情下的市场适配性，识别极端事件全周期模型核心参数的时变规律，为高波动市场环境下的衍生品定价与风险管理提供理论支撑与实证依据。

\section{研究目的与主要内容}

针对上述核心研究问题，本文以 3/2 随机波动率模型下的欧式期权定价为核心研究对象，通过理论构建、数值实验与实证分析三个层面的系统研究，全面评估不同定价方法的全场景性能，厘清其适用边界，同时验证 3/2 模型在极端市场环境下的实证适配性。本文的核心研究目的与主要研究内容分为以下三个方面：

第一，首次系统推导并实现逆高斯（Inverse Gaussian, IG）近似精确模拟、Gamma 近似精确模拟、Poisson-Gamma 近似精确步进、拟指数（Quasi-Exponential, QE）近似精确步进方法，以及傅里叶余弦（COS）定价方法在 3/2 纯扩散模型下的欧式期权定价方案，构建 3/2 模型下欧式期权定价的完整方法体系。本文将十一类主流定价方法划分为时间步进类条件蒙特卡洛方法、终端精确与近似精确模拟方法、傅里叶频域定价方法三大范式，系统推导每类方法的核心理论逻辑、数值实现流程、误差特性与收敛性质，解决现有研究中方法体系分散、理论边界不清晰、核心近似精确方法在 3/2 模型中应用空白的问题，为后续的数值对比与实证检验构建统一的理论框架与可复现的实现方案。

第二，通过多场景数值实验，完成不同定价方法的系统性性能评价。本文设计 7 组覆盖常规市场与极端参数环境的测试案例，以数值稳定性、定价精度、计算效率为核心评价指标，重点对标本文首次引入 3/2 模型的 IG、Gamma、Poisson-Gamma、QE 近似精确方法与 COS 方法，和传统 Euler/Milstein 离散格式、终端精确模拟方法的性能差异，量化对比十一类方法在不同参数组合下的表现，明确各类方法的优势场景与核心缺陷，形成清晰的方法性能分层体系，解决现有研究中方法性能评价片面、极端场景鲁棒性检验不足的问题，为不同应用场景下的方法选择提供量化依据。

第三，基于真实极端市场事件，完成 3/2 模型的实证校准与定价方法的极端场景再检验。本文以 2025 年 10 月比特币市场闪崩极端事件为研究样本，构建 “长期固定参数时间序列估计 - 时变参数横截面校准” 的两步实证框架，识别崩盘事件全周期 3/2 模型核心参数的时变规律，验证模型对极端行情下波动率微笑演化的拟合能力；同时基于校准得到的真实市场参数，构建极端场景测试案例，完成定价方法的性能再检验，重点验证本文首次引入的各类近似精确方法与 COS 方法在真实极端市场中的鲁棒性，并与学术基准算例结果进行横向对比，厘清理想化环境与真实市场中方法性能的核心差异，最终形成极端市场环境下的方法选择指引，弥补现有研究中理论算例与真实市场脱节的不足。

\section{论文结构安排}

本文共设置七个章节，遵循“理论基础-方法构建-数值实验-实证检验-结论展望”的逻辑主线展开，各章节的核心内容与结构安排如下：

第一章为引论。本章系统阐述本文的研究背景与研究意义，梳理全球及国内期权定价理论的发展脉络，明确本文针对 3/2 随机波动率模型的核心研究问题、研究目的与主要研究内容，并介绍全文的章节结构安排。

第二章为文献综述。本章系统梳理随机波动率模型的理论演进脉络，精确模拟定价方法的研究现状，总结 3/2 模型在金融衍生品市场中的应用研究与国内相关研究现状，最终通过文献述评明确现有研究的不足，确立本文的学术定位与研究边际贡献。

第三章为 3/2 随机波动率模型设定。本章详细介绍 3/2 随机波动率模型的核心理论设定，推导风险中性测度下标的资产价格与瞬时方差过程的随机微分方程，推导模型积分方差的条件分布性质与标的价格的矩母函数表达，厘清模型核心参数的经济含义，为后续定价方法的构建与实证校准奠定理论基础。

第四章为期权定价方法与算法实现。本章在 3/2 纯扩散模型的统一理论框架下，系统构建三大范式十一类期权定价方法的完整实现体系，核心创新在于首次完成了逆高斯（IG）、Gamma 近似精确模拟、Poisson-Gamma 近似精确步进、QE 近似精确步进四类近似精确方法，以及 COS 定价方法在 3/2 模型下的系统性理论推导与数值实现；分别推导各类方法的核心理论公式，给出每类方法的标准化数值实现流程，并对其误差特性、收敛性质与适用场景进行理论分析，为后续的数值实验提供统一的方法基准。

第五章为数值实验设定与结果对比分析。本章设定统一的数值实验环境与评价指标体系，设计 7 组覆盖常规市场与极端参数环境的测试案例，系统完成十一类定价方法的数值测试，量化对比不同方法的数值稳定性、定价精度与计算效率，最终形成方法性能的分层体系与场景化选择建议。

第六章为 3/2 随机波动率模型在加密货币极端行情中的实证校准与定价检验。本章以 2025 年 10 月比特币市场闪崩等极端事件为研究样本，构建“长期固定参数时间序列估计-时变参数横截面校准-定价方法极端场景再检验”的完整实证框架，识别崩盘事件全周期模型核心参数的时变规律与模型拟合效果，验证模型对极端行情下波动率微笑演化的拟合能力；同时基于真实市场校准参数完成定价方法的极端场景性能再检验，与第五章的学术基准算例结果进行横向对比，明确极端市场环境下不同定价方法的适用边界，重点验证本文引入的各类近似精确方法与 COS 方法在真实极端市场中的工程适用性。

第七章为结论与展望。本章系统总结全文的核心研究结论，客观分析本研究存在的局限性，并从模型扩展、方法优化、实证范围拓展、风险管理应用等维度，对未来相关研究方向进行展望。